% !TEX program = xelatex
% Template ABNT para resenhas críticas derivado do arquivo resenha_walden2.tex

\documentclass[12pt, a4paper, oneside]{abntex2}

% --- PACOTES PRINCIPAIS ---
\usepackage{fontspec}
\usepackage[brazil]{babel}
\usepackage[alf,abnt-show-options=none,abnt-emphasize=bf]{abntex2cite}
\usepackage{geometry}
\usepackage{hyperref}
\usepackage{etoolbox}
\usepackage{xparse}
\usepackage{indentfirst}
\usepackage{graphicx}
\usepackage{titlesec}
\AtBeginDocument{\let\tt\ttfamily}

\geometry{left=3cm, top=3cm, right=2cm, bottom=2cm}
\setmainfont{Arial}

% --- PARÁGRAFOS E ESPAÇAMENTO ---
\setlength{\parindent}{1.25cm}
\setlength{\parskip}{0pt}
\OnehalfSpacing
\raggedbottom

% --- ESPAÇAMENTO ENTRE TÍTULOS E TEXTO ---
\titlespacing*{\chapter}{0pt}{0pt}{1.5\baselineskip}
\titlespacing*{\section}{0pt}{1.5\baselineskip}{1.5\baselineskip}
\titlespacing*{\subsection}{0pt}{1.5\baselineskip}{1.5\baselineskip}
\titlespacing*{\subsubsection}{0pt}{1.5\baselineskip}{1.5\baselineskip}

% --- Nomes das seções pré-textuais ---
\renewcommand{\resumoname}{RESUMO}
\renewcommand{\abstractname}{ABSTRACT}
\renewcommand{\contentsname}{\textbf{\MakeUppercase{SUMÁRIO}}}

% --- Ajustes de fonte para títulos ---
\renewcommand{\ABNTEXchapterfont}{\bfseries\normalsize}
\renewcommand{\ABNTEXchapterfontsize}{\normalsize}
\renewcommand{\ABNTEXsectionfont}{\bfseries\normalsize}
\renewcommand{\ABNTEXsectionfontsize}{\normalsize}
\renewcommand{\ABNTEXsubsectionfont}{\normalfont\normalsize}
\renewcommand{\ABNTEXsubsectionfontsize}{\normalsize}
\setboolean{ABNTEXupperchapter}{true}
\setboolean{ABNTEXuppersection}{false}
\setboolean{ABNTEXuppersubsection}{false}

% --- Numeração de página ---
\makepagestyle{padraopagina}
\makeevenhead{padraopagina}{}{}{\fontsize{10}{12}\selectfont\thepage}
\makeoddhead{padraopagina}{}{}{\fontsize{10}{12}\selectfont\thepage}
\makeevenfoot{padraopagina}{}{}{}
\makeoddfoot{padraopagina}{}{}{}

% --- Citações diretas longas ---
\renewenvironment{citacao}{%
  \list{}{%
    \listparindent=\parindent
    \parsep=0pt plus 1pt
    \leftmargin=4cm
    \rightmargin=0cm
  }%
  \item\relax\fontsize{10}{15}\selectfont
}{\endlist}

% --- DADOS EDITÁVEIS DO TRABALHO ---
\newcommand{\nomeinstituicao}{Universidade do Estado da Bahia -- UNEB}
\newcommand{\nomedepartamento}{Departamento de Ciências da Vida -- DCV}
\newcommand{\nomecurso}{Curso de Bacharelado em Psicologia}
\newcommand{\nomeautor}{Lucas Camilo Carvalho}
\newcommand{\nomecidade}{Salvador}
\newcommand{\anoacademico}{2026}
\newcommand{\nomedisciplina}{Teorias da Personalidade}
\newcommand{\nomeorientador}{Cláudia}
\newcommand{\tituloprincipal}{SÍNTESE REFLEXIVA}
\newcommand{\titulocomplemento}{Da queda livre ao encontro com o outro nas redes sociais: um estudo do narcisismo}
\newcommand{\tipodocumento}{Síntese reflexiva}
\newcommand{\logoinstituicao}{logo_uneb.png}

\newcommand{\formatartitulo}{%
  \ifx\titulocomplemento\empty
    {\bfseries\normalsize\MakeUppercase{\tituloprincipal}}%
  \else
    {\bfseries\normalsize\MakeUppercase{\tituloprincipal}: \normalfont\normalsize\titulocomplemento}%
  \fi
}
\titulo{\formatartitulo}

% --- METADADOS ---
\autor{\textbf{\MakeUppercase{\nomeautor}}}
\orientador{}
\coorientador{}
\instituicao{%
  \textbf{\MakeUppercase{\nomeinstituicao}}\\
  \textbf{\MakeUppercase{\nomedepartamento}}\\
  \textbf{\MakeUppercase{\nomecurso}}
}
\tipotrabalho{\textbf{\MakeUppercase{\tipodocumento}}}
\local{\textbf{\nomecidade}}
\data{\textbf{\anoacademico}}

% --- DESCRIÇÃO NA FOLHA DE ROSTO ---
\newcommand{\folhaderostodescricao}{%
  Trabalho apresentado como requisito parcial para a obtenção de nota na disciplina \nomedisciplina, sob orientação de \nomeorientador, no \nomecurso\ da \nomeinstituicao.
}
\preambulo{}

\newcommand{\imprimircapaformatada}{%
  \begin{capa}
    \begin{center}
      \vspace*{-1.73cm}
      \IfFileExists{\logoinstituicao}{\includegraphics[width=5cm]{\logoinstituicao}\par}{}
      \vspace{1cm}
      {\bfseries\normalsize\MakeUppercase{\nomeinstituicao}\par}
      {\bfseries\normalsize\MakeUppercase{\nomedepartamento}\par}
      {\bfseries\normalsize\MakeUppercase{\nomecurso}\par}
      \vspace{3cm}
      {\bfseries\normalsize\MakeUppercase{\nomeautor}\par}
      \vspace{3cm}
      {\formatartitulo\par}
      \vfill
      {\bfseries\normalsize\nomecidade\par}
      {\bfseries\normalsize\anoacademico\par}
    \end{center}
  \end{capa}
}

\newcommand{\imprimirfolhaderostoformatada}{%
  \begin{folhaderosto}
    \SingleSpacing
    {\centering\bfseries\normalsize\MakeUppercase{\nomeautor}\par}
    \vspace{\fill}
    {\centering\formatartitulo\par}
    \vspace{1.5cm}
    {
      \SingleSpacing
      \leftskip=7cm
      \parindent=0pt
      \noindent\folhaderostodescricao\par
    }
    \vspace{\fill}
    {\centering\bfseries\normalsize\nomecidade\par}
    {\centering\bfseries\normalsize\anoacademico\par}
  \end{folhaderosto}
}

% --- Configurações de seções ---
\setsecheadstyle{\normalsize\bfseries}
\setsubsecheadstyle{\normalsize}
\setsubsubsecheadstyle{\normalsize}
\setlength{\afterchapskip}{1.5\baselineskip}
\setaftersecskip{1.5\baselineskip}
\setaftersubsecskip{1.5\baselineskip}
\setaftersubsubsecskip{1.5\baselineskip}
\setlength{\cftbeforechapterskip}{0pt}
\setlength{\cftbeforesectionskip}{0pt}
\setlength{\cftbeforesubsectionskip}{0pt}
\cftsetindents{chapter}{0em}{1em}
\cftsetindents{section}{0em}{2em}
\cftsetindents{subsection}{0em}{2.5em}
\cftsetindents{subsubsection}{0em}{3em}
\setlength{\cftlastnumwidth}{1em}
\renewcommand{\cftchapterfont}{\bfseries\MakeUppercase}
\renewcommand{\cftchapterpagefont}{\bfseries}
\renewcommand{\cftchapteraftersnum}{\space}
\renewcommand{\chapternumberline}[1]{#1\hspace{0.3em}}
\newcommand{\overridetoclengths}{%
  \let\oldsetlength\setlength
  \def\setlength##1##2{%
    \ifx##1\cftchapterindent
      \oldsetlength{\cftchapterindent}{0pt}%
    \else\ifx##1\cftchapternumwidth
      \oldsetlength{\cftchapternumwidth}{1em}%
    \else
      \oldsetlength{##1}{##2}%
    \fi\fi}%
}
\newcommand{\restoretoclengths}{\let\setlength\oldsetlength}

% --- Chaves para exibição ---
\newif\ifmostrarcapa
\newif\ifmostrarfolhaderosto
\newif\ifmostrarresumo
\newif\ifmostrarabstract
\newif\ifmostrarsumario
\newif\ifmostrartexto
\mostrarcapatrue
\mostrarfolhaderostotrue
\mostrarresumotrue
\mostrarabstractfalse
\mostrarsumariotrue
\mostrartextotrue
\renewcommand{\cftsectionfont}{\bfseries}
\renewcommand{\cftsectionpagefont}{\bfseries}
\renewcommand{\cftsubsectionfont}{\normalfont}
\renewcommand{\cftsubsectionpagefont}{\normalfont}
\renewcommand{\cftsubsubsectionfont}{\normalfont}
\renewcommand{\cftsubsubsectionpagefont}{\normalfont}
\renewcommand{\cftdotsep}{1}

% --- Referências ---
\renewcommand{\bibname}{\textbf{\MakeUppercase{REFERÊNCIAS}}}
\providecommand{\autorcite}{}
\RenewDocumentCommand{\autorcite}{s O{} m}{%
  \IfBooleanTF{#1}
    {\hyperlink{cite.#3}{\citeyear{#3}\IfBlankTF{#2}{}{, p.~#2}}}
    {\hyperlink{cite.#3}{(\citeauthoronline{#3}, \citeyear{#3}\IfBlankTF{#2}{}{, p.~#2})}}}

% --- Biblioteca BibTeX de exemplo ---
\begin{filecontents*}{\jobname.bib}
@article{pizzimenti2019,
  author    = {Enzo Cléto Pizzimenti and Isis Graziele Da Silva and Ivan Ramos Estêvão},
  title     = {Da queda livre ao encontro com o outro nas redes sociais: um estudo do narcisismo},
  journal   = {Trivium: Estudos Interdisciplinares},
  year      = {2019},
  volume    = {Ano XI},
  number    = {Ed. 1},
  pages     = {85--98}
}
\end{filecontents*}

\begin{filecontents*}{\jobname.toc}
\changetocdepth {4}
\babel@toc {brazil}{}\relax
\end{filecontents*}

\begin{document}
\addtocontents{toc}{\protect\setlength{\protect\cftchapterindent}{0pt}}
\addtocontents{toc}{\protect\setlength{\protect\cftchapternumwidth}{1em}}

% --- PRETEXTUAL ---
\pretextual
\pagestyle{empty}
\SingleSpacing
\ifmostrarcapa
  \imprimircapaformatada
\fi
\ifmostrarfolhaderosto
  \imprimirfolhaderostoformatada
\fi
\ifmostrarresumo
  \begingroup
  \SingleSpacing
  \setlength{\parindent}{0pt}
  \begin{resumo}
  Esta síntese avalia o artigo ``Da queda livre ao encontro com o outro nas redes sociais: um estudo do narcisismo'', de Pizzimenti, Silva e Estêvão (2019). O texto mostra como as pessoas usam o termo narcisismo hoje. Ele estuda os efeitos disso na mente atual. Para isso, usa o episódio ``Queda Livre'' da série Black Mirror. A obra aponta a tensão entre exibir a imagem própria como defesa e a chance de encontros reais no meio digital.
  \vspace{\baselineskip}
  \textbf{Palavras-chave:} Narcisismo. Redes sociais. Psicanálise. Pós-modernidade.
  \end{resumo}
  \endgroup
\fi
\ifmostrarabstract
  \begingroup
  \SingleSpacing
  \setlength{\parindent}{0pt}
  \begin{resumo}[ABSTRACT]
  \begin{otherlanguage*}{english}
  This reflective synthesis analyzes the article ``From the free fall to the encounter with the other in the social networks: a narcissism's study'', by Pizzimenti, Silva, and Estêvão (2019). The text discusses the appropriation of the term narcissism by popular discourse and investigates its effects on contemporary subjectivity, using the Black Mirror episode ``Nosedive'' as an analytical key. The synthesis highlights the tension between the spectacle of self-image as a defense against postmodern helplessness and the possibilities of genuine encounter through alterity on digital platforms.
  \vspace{\baselineskip}
  \textbf{Keywords:} Narcissism. Social networks. Psychoanalysis. Postmodernity.
  \end{otherlanguage*}
  \end{resumo}
  \endgroup
\fi
\ifmostrarsumario
  \cleardoublepage
  {
    \OnehalfSpacing
    \overridetoclengths
    \pdfbookmark{\contentsname}{toc}
    \tableofcontents*
    \restoretoclengths
  }
  \cleardoublepage
\fi

% --- TEXTUAL ---
\ifmostrartexto
\OnehalfSpacing
\textual
\pagestyle{padraopagina}
\aliaspagestyle{chapter}{padraopagina}

\chapter{Introdução}
Os autores fazem uma busca na visão da psicanálise sobre o uso de redes sociais\nocite{pizzimenti2019}. O senso comum adotou a palavra narcisismo para falar de vaidade ou foco em si mesmo. Mas o termo tem um peso maior na visão de Freud.

Para entender isso, o artigo usa a série \textit{Black Mirror}. No episódio ``Queda Livre'', a sociedade julga as pessoas por notas públicas. Isso funciona como um aviso sobre a lógica das redes de hoje, como o Instagram. O texto reflete se cuidar da autoimagem afasta o sujeito ou cria pontes para o outro.

\chapter{A espetacularização do eu e a demanda por avaliação}
Na trama, Lacie Pound pauta sua vida pela nota dos outros. Sua nota dita seu valor social e o seu acesso a bens e locais. A ficção copia a vida real nas redes atuais. Nelas, a vida privada vira um show e um produto.

Os autores citam Zygmunt Bauman, Christopher Lasch e Gilles Lipovetsky para situar o leitor. A época atual perdeu as garantias fixas do passado, como a família e o Estado. Com isso, o sujeito vive em um mundo de trocas rápidas e dúvidas. Sem apoios firmes, a pessoa tenta criar quem ela é a todo tempo. Ela passa a buscar sempre a aprovação de fora.

É nesse ponto que o narcisismo ganha peso. A busca por likes não é apenas um traço de ego alto. Ela age como defesa contra a dor de viver num mundo sem laços fortes. Quando a heroína da série tenta subir sua nota ao lado de pessoas ricas, nota-se essa defesa em ação. Ela tenta forjar um eu ideal a partir da visão do outro. O fim disso é tentar rever a ilusão de completude humana.

\chapter{O narcisismo: entre a defesa e o encontro}
A visão da teoria resgata Freud. Ele ensina que o narcisismo não é um traço ruim ao acaso. Ele é uma fase normal da vida. O sujeito forma a si mesmo ao tomar o próprio corpo como alvo de amor. Depois, lança para o futuro o sonho de rever essa união total.

O grande risco, exposto na série, é focar apenas na imagem. Isso pode virar um ciclo sem fim. Assim, o sujeito exclui o outro real. Se quem publica não liga para \textit{quem} curte, mas apenas para o total de likes, o laço é vazio. Ele não acha o outro ali. Ele acha apenas seu reflexo num número frio. Isso leva ao tédio e ao vazio profundo que Lasch descreve.

Ainda assim, o uso das redes pode ter um lado bom. Os autores mostram que a rede também pode ser o palco de um encontro real. Quando a postagem se volta para alguém, a troca ganha vida. O lugar digital pode ajudar o sujeito a se firmar sem seguir a massa.

\chapter{Considerações finais}
Ler este artigo traz uma visão nova sobre as redes sociais. O texto afasta o leitor de apenas julgar o excesso de narcisismo de forma moral. A busca por atenção surge como sintoma do tempo atual solto e sem bases.

No fim da série, Lacie perde todos os pontos e vai presa. De forma irônica, a cela a liberta para gritar sua dor real. Essa cena alerta sobre a prisão de tentar ser perfeito o tempo todo.

Para os estudantes de psicologia, o texto dá boas bases. Ajuda a ouvir a dor de quem vive refém das telas. Ouvir a queixa de quem não se sente visto requer ver a tela não como o mal. A tela é só o novo palco dos velhos dramas da falta humana.
\fi

% --- PÓS-TEXTUAL ---
\postextual
\addtocontents{toc}{\protect\setlength{\protect\cftchapterindent}{0pt}}
\addtocontents{toc}{\protect\setlength{\protect\cftchapternumwidth}{3.2em}}
\bibliographystyle{abntex2-alf-local}
\renewcommand{\bibname}{\MakeUppercase{REFERÊNCIAS}}
\bibliography{\jobname}

\end{document}
